% acknowledgements

\begingroup
\setlength{\parindent}{0pt}
\setlength{\parskip}{\baselineskip}

Desidero esprimere il mio più sincero ringraziamento a tutte le persone che, ciascuna a modo proprio, hanno reso possibile il raggiungimento di questo traguardo.

Il mio primo pensiero va al professor \textbf{Paolo Garza}, relatore di questa tesi. La sua guida, la disponibilità e il rigore scientifico sono stati un punto di riferimento costante in ogni fase del lavoro.

Un ringraziamento sentito va al mio tutor aziendale \textbf{Marco Bartolini}, per avermi accolto in Data Reply e per la fiducia che ha riposto in me fin dall'inizio. Il confronto con lui ha arricchito ogni fase del lavoro, regalandomi una prospettiva concreta sul mondo professionale.

Un pensiero grato va al \textbf{Politecnico di Torino}, che in questi anni mi ha offerto una formazione degna di uno dei migliori atenei in Italia se non nel mondo. Non è stato un percorso semplice: le estreme difficoltà, le ansie, il peso di ogni singolo esame e le aspettative, mie e di chi mi circondava, mi hanno messo duramente alla prova. È stato proprio questo carico, però, a temprarmi e a forgiare il me di adesso, insegnandomi una resilienza che non credevo di avere. Più dei singoli esami, porterò con me gli anni vissuti qui: le persone incontrate, le sfide affrontate e la consapevolezza di essere cresciuto, dentro e fuori dall'aula.

Porto nel cuore gli anni trascorsi in Unikore. Un grazie di vero affetto va ai miei ex colleghi, i ``Gentlemen'' --- come amiamo chiamarci tra noi --- \textbf{Vito}, \textbf{Zaf}, \textbf{Gab} e \textbf{Filippo}: con loro ho condiviso non solo il lavoro, ma risate, difficoltà e una complicità che il tempo non ha scalfito, e che ancora oggi, nonostante le strade diverse, sento intatta. Un pensiero speciale va anche a \textbf{Giovanni}, ex collega e amico, che nonostante i percorsi diversi ha sempre trovato del tempo da dedicarmi e dato modo di esprimermi liberamente da vero amico.

Ringrazio i miei cari amici del liceo, \textbf{Filippo}, \textbf{Giuseppe} e \textbf{Francesco}: amicizie nate sui banchi di scuola e cresciute negli anni, diventate punti fermi a cui tornare sempre con la stessa spontaneità e la stessa ironia di un tempo. La vicinanza che mi avete dato, anche a distanza, vale più di quanto riesca a dire.

Un grazie sincero al mio amico \textbf{Christian} e alla sua fidanzata \textbf{Elsa}, per le serate, le avventure e quella leggerezza che loro tra tanti hanno saputo regalarmi: grazie per ricordarmi quanto sia importante non prendersi mai troppo sul serio e il valore di un'amicizia costante.

Un pensiero affettuoso va alla mia amica di infanzia \textbf{Viviana}, che da poco è salita anche lei a Torino: ritrovarci in una città nuova, dopo esserci conosciuti da bambini, ha reso questo posto ancora più familiare e mi ha regalato una complicità fatta di ricordi condivisi e di un affetto che la distanza non ha mai scalfito.

Un grazie anche ai ragazzi della residenza universitaria \textbf{Borsellino}, con cui ho condiviso la vita da studente fuorisede: le cene di corsa, le serate in sala studio e nel cortile, e quel senso di comunità che, fin dal primo giorno, ha fatto sentire Torino come casa e non mi ha mai fatto sentire solo.

Un grazie anche a \textbf{Pietro} e \textbf{Carmine}, conosciuti qui a Torino, per le piacevoli serate e la bella compagnia condivise: momenti spensierati che hanno reso più leggera la vita fuori dai libri, tra un episodio di House of the Dragon e dibattiti sulla qualità cinematografica di Obsession.

Un posto a parte meritano i miei insostituibili amici conosciuti qui a Torino --- \textbf{Giada}, \textbf{Samuele} e \textbf{Nicolò} --- persone che ormai considero dei fratelli e a cui voglio un bene dell'anima. Avete trasformato una città in cui ero arrivato da solo in un posto che sento mio: con voi ho diviso giorni difficili e serate indimenticabili, dubbi e risate fino a tarda notte, battaglie a Mario Party e prese in giro a Coco Rido. Mi avete aiutato ogni volta che ne ho avuto bisogno e non vi siete mai fatti indietro nemmeno un po'. Avervi trovati lungo il cammino è uno dei regali più grandi di questi anni, e rimanere qui a Torino non sarebbe la stessa senza di voi.

Devo un affetto profondo ai miei amici di infanzia \textbf{Marco}, \textbf{Giorgia} e i loro genitori \textbf{Giovanni} e \textbf{Silvana}: una seconda famiglia fin da bambino, con cui ho condiviso momenti indimenticabili e che, nonostante il tempo e la distanza, continua a rappresentare un porto sicuro in cui so di poter sempre tornare.

Un ringraziamento di cuore va a tutta la mia famiglia da parte di padre, i \textbf{Lopez}. Sono così numerosi che non posso nominarli uno ad uno, eppure nessuno di loro mi è estraneo: agli zii, ai cugini e a tutti i parenti che hanno sempre fatto sentire la loro presenza, il loro calore e il loro orgoglio per me, va il mio affetto e la mia riconoscenza, raccolti in un solo grande abbraccio. Un pensiero particolare va poi a chi purtroppo non c'è più: anche se non siete qui a condividere questo traguardo, è come se ci foste ancora, perché continuate a vivere nei ricordi, in tutto ciò che avete trasmesso e attraverso mio padre.

Con affetto ringrazio i miei suoceri \textbf{Enza} e \textbf{Roberto} e mia cognata \textbf{Chiara}, che mi hanno accolto come un figlio e come un fratello, facendomi sentire parte di una famiglia nuova con una naturalezza che non do mai per scontata, e insieme a loro tutta la loro grande famiglia di zii, parenti e cugini, per il calore genuino con cui mi hanno sempre circondato. Un pensiero riconoscente va ai miei zii \textbf{Teresa} e \textbf{Dino}, presenza costante e affettuosa che ha seguito passo dopo passo il mio percorso, e ai miei cugini \textbf{Andrea} e \textbf{Riccardo}, con cui ho condiviso, tappa dopo tappa, la crescita e innumerevoli ricordi.

Ringrazio i miei nonni \textbf{Giuseppe} e \textbf{Cettina}, che con la loro dolcezza e i loro insegnamenti hanno riempito la mia infanzia di affetto e di sicurezza. Li ammiro per l'esempio silenzioso di vita che continuano a rappresentare: la loro sobrietà, la loro laboriosità e l'amore generoso che non mi hanno mai lesinato sono valori che porto nel cuore e radici da cui attingo ancora oggi. Essere loro nipote è uno dei privilegi più grandi della mia vita, e sapere che sono orgogliosi di me rende ogni traguardo ancora più bello.

Alla mia fidanzata \textbf{Giulia} va il mio grazie più profondo. La tua presenza, la tua pazienza e il tuo amore mi hanno sostenuto nei momenti in cui la stanchezza aveva la meglio e reso più belli quelli felici; hai creduto in me anche quando stentavo a farlo io, e sai celebrare le mie piccole vittorie come se fossero tue. Negli otto anni passati insieme abbiamo attraversato anche momenti difficili, ma essere riusciti ad arrivare fin qui, dopo tutto questo tempo, è la prova concreta di quanto sia solido e profondo il nostro rapporto. Senza di te, ogni tappa di questo percorso sarebbe stata molto più faticosa e decisamente meno felice. Non vedo l'ora di iniziare la nostra prossima avventura, di condividere una casa e di costruire insieme il nostro futuro qui a Torino: è con te che voglio affrontare tutto ciò che verrà.

Infine, e con tutto il cuore, ringrazio i miei \textbf{genitori}. Il vostro amore incondizionato, la vostra fiducia in me e i sacrifici che avete affrontato, spesso in silenzio, per darmi ogni opportunità sono il fondamento di tutto ciò che ho raggiunto. Avete sempre creduto nelle mie capacità, anche quando io per primo ne dubitavo, e mi avete sostenuto in ogni scelta, lasciandomi libero di sbagliare e di crescere. Avete sempre fatto di tutto per non farmi mai mancare nulla, anche nei momenti più delicati: lo ricordo bene in questi ultimi mesi. Prima che la recruiter di Intesa Sanpaolo mi contattasse per il ruolo vero e proprio era già trascorso un mese e mezzo, e pur di coronare il sogno di entrare in una multinazionale così grande ero disposto ad accettare anche sei mesi di tirocinio, nella speranza di aumentare le mie probabilità di assunzione, ben sapendo che avrebbe significato ``sopravvivere'' per altri sei mesi senza la stabilità di un lavoro. In quella situazione non mi avete mai fatto pesare nulla: mi avete incoraggiato a fare la scelta che io stesso ritenessi più giusta, assicurandomi che avreste pensato voi a mandarmi avanti e a non farmi mancare niente. E questo, del resto, è solo uno dei tanti esempi che porto nel cuore: in ogni momento di dubbio ci siete sempre stati, pronti a sostenermi senza chiedere nulla in cambio. Siete il mio esempio di vita, il mio porto sicuro e la ragione per cui ho potuto coltivare le mie ambizioni lontano da casa. Nessun traguardo mi appartiene davvero, se non lo condivido con voi. A voi dedico questa tesi, perché senza di voi nulla di tutto questo sarebbe stato possibile.

Un pensiero speciale, accanto ai miei genitori, va anche ai miei piccoli compagni a quattro zampe: al nostro chihuahua \textbf{Toby}, che con le sue feste e la sua ingombrante dolcezza rende ogni ritorno a casa più caldo. E un pensiero pieno di tenerezza al suo predecessore, \textbf{Toby Senior}, venuto a mancare ormai da diversi anni: non passa giorno senza che pensi a lui, e mi piace credere che continui a vivere, in qualche modo, nel piccolo Toby di oggi --- stesso nome, stesso affetto immenso, come se non se ne fosse mai andato del tutto.

Un'ultima considerazione, dedicata a tutti voi: scelgo di non aggiungere i classici ringraziamenti a me stesso, perché credo che ognuno di voi abbia contribuito a questo traguardo e, soprattutto, a definire chi sono oggi. Ogni incontro, ogni consiglio e ogni gesto di affetto hanno lasciato una traccia in me, e questo risultato è tanto vostro quanto mio.

Questa tesi chiude un capitolo e ne apre uno nuovo: il mio ingresso nel mondo del lavoro. Porto con me tutto ciò che ho imparato e, soprattutto, le persone nominate in queste pagine, mentre affronto con entusiasmo la mia prima tappa professionale come \emph{Machine Learning Engineer} presso il grattacielo Intesa Sanpaolo. Lavorare in alto, sospesi tra i vetri di quella torre, mi regalerà ogni giorno uno sguardo che abbraccia Torino intera: mi ricorderà da dove sono partito e quanto strada ho fatto. Perché se è vero che si arriva in alto, è altrettanto vero che non si sale mai da soli. Quel panorama, come ogni traguardo di queste pagine, è anche vostro. Grazie.

\endgroup
\vspace*{5\baselineskip}
